\documentclass[11pt,a4paper]{article}

\usepackage[margin=2.5cm]{geometry}
\usepackage{amsmath, amssymb, amsthm}
\usepackage{bm}
\usepackage{enumitem}
\usepackage{booktabs}
\usepackage{parskip}

\newtheorem{definition}{Definition}
\newtheorem{proposition}{Proposition}

\title{\textbf{Pity Counter Collision Probability}\\[0.3em]
\large A Markov Chain Analysis of Simultaneous Pity Events}
\author{}
\date{}

\begin{document}
\maketitle

\section{Overview}

The loot system contains two independent pity mechanics. The S-tier pity counter forces an S-tier outcome after 80 consecutive non-S-tier rolls. The A-tier pity counter forces an A-tier outcome after 10 consecutive non-A-tier rolls. When both counters reach their maximum simultaneously, the system must resolve a collision. We derive the long-run probability of this event.

\section{Model}

\begin{definition}
Let the \textbf{S-tier pity counter} be modelled as a discrete-time Markov chain with state space $\mathcal{S}_S = \{0, 1, 2, \ldots, 79\}$ and transition probabilities:
\begin{align}
P_S(i, 0)   &= p_S           &&\text{for } i \in \{0, 1, \ldots, 78\} \quad \text{(natural hit, reset)} \\
P_S(i, i+1) &= 1 - p_S       &&\text{for } i \in \{0, 1, \ldots, 78\} \quad \text{(miss, advance)} \\
P_S(79, 0)  &= 1             &&\text{(pity, forced reset)}
\end{align}
where $p_S = 0.005$.
\end{definition}

\begin{definition}
Let the \textbf{A-tier pity counter} be modelled as a discrete-time Markov chain with state space $\mathcal{S}_A = \{0, 1, 2, \ldots, 9\}$ and identical transition structure with $p_A = 0.10$.
\end{definition}

The transition matrix for each chain has the form:
\[
\mathbf{P} = 
\begin{pmatrix}
p     & 1-p   & 0     & \cdots & 0     & 0     \\
p     & 0     & 1-p   & \cdots & 0     & 0     \\
p     & 0     & 0     & \cdots & 0     & 0     \\
\vdots& \vdots& \vdots& \ddots & \vdots& \vdots\\
p     & 0     & 0     & \cdots & 0     & 1-p   \\
1     & 0     & 0     & \cdots & 0     & 0     
\end{pmatrix}
\]

where each row sums to 1 and the final row reflects the forced pity reset.

\section{Stationary Distribution}

\begin{proposition}
For a pity counter with natural hit probability $p$ and $n$ states $\{0, 1, \ldots, n-1\}$, the stationary distribution is:
\[
\pi(k) = \frac{(1-p)^k \cdot p}{1 - (1-p)^n} \qquad \text{for } k = 0, 1, \ldots, n-1
\]
\end{proposition}

\begin{proof}
State $k$ is reached only by $k$ consecutive misses from state 0, giving the relationship:
\[
\pi(k) = (1-p)^k \cdot \pi(0)
\]

Enumerating every state:
\begin{alignat*}{2}
\pi(0)   &= (1-p)^0 \cdot \pi(0)   &&= 1 \cdot \pi(0) \\
\pi(1)   &= (1-p)^1 \cdot \pi(0)   && \\
\pi(2)   &= (1-p)^2 \cdot \pi(0)   && \\
         &\;\;\vdots                 && \\
\pi(n-1) &= (1-p)^{n-1} \cdot \pi(0) &&
\end{alignat*}

Every state is now expressed in terms of the single unknown $\pi(0)$. To solve for it, we apply the normalisation constraint --- the probability of being in \emph{some} state must equal 1:
\[
\sum_{k=0}^{n-1} \pi(k) = 1
\]

Substituting each $\pi(k)$:
\begin{align}
\pi(0) + \pi(1) + \pi(2) + \cdots + \pi(n-1) &= 1 \notag \\[6pt]
(1-p)^0 \cdot \pi(0) + (1-p)^1 \cdot \pi(0) + \cdots + (1-p)^{n-1} \cdot \pi(0) &= 1 \notag \\[6pt]
\bigl((1-p)^0 + (1-p)^1 + \cdots + (1-p)^{n-1}\bigr) \cdot \pi(0) &= 1 \label{eq:factored}
\end{align}

The bracketed sum is a geometric series with first term $a = 1$, common ratio $r = (1-p)$, and $n$ terms:
\[
\sum_{k=0}^{n-1} (1-p)^k = \frac{1 - (1-p)^n}{1 - (1-p)} = \frac{1 - (1-p)^n}{p}
\]

Substituting this result back into equation~\eqref{eq:factored}:
\[
\frac{1 - (1-p)^n}{p} \cdot \pi(0) = 1
\]

Solving for $\pi(0)$:
\[
\pi(0) = \frac{p}{1 - (1-p)^n}
\]

Substituting back into $\pi(k) = (1-p)^k \cdot \pi(0)$:
\[
\pi(k) = \frac{(1-p)^k \cdot p}{1 - (1-p)^n}
\]
\end{proof}

\section{Application}

\subsection{S-tier pity state}

For the S-tier chain ($p_S = 0.005$, $n = 80$), the probability of being in the maximum pity state is:
\[
\pi_S(79) = \frac{0.995^{79} \times 0.005}{1 - 0.995^{80}} = \frac{0.6730 \times 0.005}{1 - 0.6696} = \frac{0.003365}{0.3304} \approx 0.01018
\]

\subsection{A-tier pity state}

For the A-tier chain ($p_A = 0.10$, $n = 10$):
\[
\pi_A(9) = \frac{0.9^{9} \times 0.1}{1 - 0.9^{10}} = \frac{0.3874 \times 0.1}{1 - 0.3487} = \frac{0.03874}{0.6513} \approx 0.05953
\]

\subsection{Collision probability}

The two chains evolve independently: the S-tier counter's transitions do not depend on the A-tier counter's state, and vice versa. Therefore, by the multiplication rule for independent events:
\[
\boxed{P(\text{collision}) = \pi_S(79) \times \pi_A(9) = 0.01018 \times 0.05953 \approx 0.000606}
\]

This corresponds to approximately a $0.06\%$ chance on any given roll.

\subsection{Expected collisions}

Over $N$ rolls, the expected number of simultaneous pity events is approximately:
\[
\mathbb{E}[\text{collisions}] \approx N \cdot P(\text{collision})
\]

\begin{table}[h]
\centering
\begin{tabular}{@{}rc@{}}
\toprule
Rolls ($N$) & Expected collisions \\
\midrule
100   & 0.06 \\
500   & 0.30 \\
1000  & 0.61 \\
2000  & 1.21 \\
3000  & 1.82 \\
\bottomrule
\end{tabular}
\end{table}

The collision is uncommon but not negligible at higher roll counts. The system resolves collisions by prioritising the rarer S-tier pity, deferring the A-tier pity to the immediately following roll.

\end{document}
